
Серов, художник, пошел на Обводной канал. Зачем он туда пошел? Покупать резину. Зачем ему резина? Чтобы сделать себе резинку. А зачем ему резинка? А чтобы ее растягивать. Вот. Что еще? А еще вот что: художник Серов поломал свои часы. Часы хорошо ходили, а он их взял и поломал. Чего еще? А более ничего. Ничего, и всё тут! И свое поганое рыло куда не надо не суй! Господи помилуй!

Жила-была старушка. Жила, жила и сгорела в печке. Туда ей и дорога!

Серов, художник, по крайней мере так рассудил.

Эх! Написать бы еще, да чернильница куда"=то исчезла.

\begin{flushright}
22~октября 1938~г.
\end{flushright}

\begin{center}
* * *
\end{center}